\documentclass{standalone}
\usepackage{circuitikz}
\usetikzlibrary{calc}
\definecolor{circuitnotemark}{HTML}{FE00FE}
\begin{document}
\begin{circuitikz}[american, line width=0.8pt]
\ctikzset{bipoles/length=1.2cm}
\ctikzset{grounds/scale=1.36}
\foreach \x in {0,2,4,6,8,10,12,14,16,18} {\foreach \y in {0,-2,-4,-6,-8,-10,-12,-14,-16,-18} {\fill[gray, opacity=0.35] (\x,\y) circle (1.1pt);}}
\node[gray, font=\scriptsize] at (0,0.84) {1};
\node[gray, font=\scriptsize] at (2,0.84) {2};
\node[gray, font=\scriptsize] at (4,0.84) {3};
\node[gray, font=\scriptsize] at (6,0.84) {4};
\node[gray, font=\scriptsize] at (8,0.84) {5};
\node[gray, font=\scriptsize] at (10,0.84) {6};
\node[gray, font=\scriptsize] at (12,0.84) {7};
\node[gray, font=\scriptsize] at (14,0.84) {8};
\node[gray, font=\scriptsize] at (16,0.84) {9};
\node[gray, font=\scriptsize] at (18,0.84) {10};
\node[gray, font=\scriptsize, anchor=east] at (-0.84,0) {a};
\node[gray, font=\scriptsize, anchor=east] at (-0.84,-2) {b};
\node[gray, font=\scriptsize, anchor=east] at (-0.84,-4) {c};
\node[gray, font=\scriptsize, anchor=east] at (-0.84,-6) {d};
\node[gray, font=\scriptsize, anchor=east] at (-0.84,-8) {e};
\node[gray, font=\scriptsize, anchor=east] at (-0.84,-10) {f};
\node[gray, font=\scriptsize, anchor=east] at (-0.84,-12) {g};
\node[gray, font=\scriptsize, anchor=east] at (-0.84,-14) {h};
\node[gray, font=\scriptsize, anchor=east] at (-0.84,-16) {i};
\node[gray, font=\scriptsize, anchor=east] at (-0.84,-18) {j};
\coordinate (x1y2) at (0,-2);
\coordinate (x5y5) at (8,-8);
\coordinate (x3y7) at (4,-12);
\coordinate (x2y6) at (2,-10);
\coordinate (x2y8) at (2,-14);
\coordinate (x2y10) at (2,-18);
\coordinate (x6y8) at (10,-14);
\coordinate (x6y10) at (10,-18);
\coordinate (x10y2) at (18,-2);
\coordinate (x10y4) at (18,-6);
\coordinate (x10y6) at (18,-10);
\coordinate (x10y8) at (18,-14);
\coordinate (x4y7) at (6,-12);
\coordinate (x9y7) at (16,-12);
\coordinate (x9y6) at (16,-10);
\coordinate (x7y2) at (12,-2);
\coordinate (x8y4) at (14,-6);
\draw (x1y2) node[ocirc]{} node[above left]{VCC}; % line 3
\node[dipchip, num pins=8, hide numbers, font=\scriptsize, /tikz/circuitikz/multipoles/dipchip/width=1.5] (part-U1) at (x5y5) {}; % line 4
\node[font=\scriptsize, anchor=north] at (part-U1.south) {$\mathrm{NE555}$}; % line 4
\node[anchor=south] at (part-U1.north) {$U_{1}$}; % line 4
\node[circuitnotemark, font=\tiny, anchor=west, xshift=2pt, inner sep=0] at (part-U1.bpin 1) {X}; % line 4
\node[circuitnotemark, font=\tiny, anchor=west, xshift=2pt, inner sep=0] at (part-U1.bpin 2) {X}; % line 4
\node[circuitnotemark, font=\tiny, anchor=west, xshift=2pt, inner sep=0] at (part-U1.bpin 3) {X}; % line 4
\node[circuitnotemark, font=\tiny, anchor=west, xshift=2pt, inner sep=0] at (part-U1.bpin 4) {X}; % line 4
\node[circuitnotemark, font=\tiny, anchor=east, xshift=-2pt, inner sep=0] at (part-U1.bpin 5) {X}; % line 4
\node[circuitnotemark, font=\tiny, anchor=east, xshift=-2pt, inner sep=0] at (part-U1.bpin 6) {X}; % line 4
\node[circuitnotemark, font=\tiny, anchor=east, xshift=-2pt, inner sep=0] at (part-U1.bpin 7) {X}; % line 4
\node[circuitnotemark, font=\tiny, anchor=east, xshift=-2pt, inner sep=0] at (part-U1.bpin 8) {X}; % line 4
\node[ground] at (x3y7) {}; % line 5
\draw (x2y6) to[R, l_=$R_{3}$, a^=$330\,\Omega$] (x2y8); % line 6
\draw (x2y8) to[leD, l_=$D_{1}$] (x2y10); % line 7
\node[anchor=west] at (2.8,-16) {$\mathrm{red}$}; % line 7
\node[ground] at (x2y10) {}; % line 8
\draw (x6y8) to[C, l_=$C_{2}$, a^=$10\,\mathrm{n}\mathrm{F}$] (x6y10); % line 9
\node[ground] at (x6y10) {}; % line 10
\draw (x10y2) to[R, l_=$R_{1}$, a^=$10\,\mathrm{k}\Omega$] (x10y4); % line 11
\draw (x10y4) to[R, l_=$R_{2}$, a^=$47\,\mathrm{k}\Omega$] (x10y6); % line 12
\draw (x10y6) to[C, l_=$C_{1}$, a^=$10\,\mu\mathrm{F}$] (x10y8); % line 13
\node[ground] at (x10y8) {}; % line 14
\draw (x1y2) -- (x10y2); % line 16
\draw (part-U1.pin 4) -| (x1y2); % line 17
\draw (part-U1.pin 1) -| (x3y7); % line 18
\draw (part-U1.pin 3) -| (x2y6); % line 19
\draw (part-U1.pin 2) -| (x4y7); % line 20
\draw (x4y7) -- (x9y7); % line 21
\draw (x9y7) -- (x9y6); % line 22
\draw (part-U1.pin 8) -| (x7y2); % line 23
\draw (part-U1.pin 7) -| (x8y4); % line 24
\draw (x8y4) -- (x10y4); % line 25
\draw (part-U1.pin 6) -| (x9y6); % line 26
\draw (x9y6) -- (x10y6); % line 27
\draw (part-U1.pin 5) -| (x6y8); % line 28
\node[circ] at (x1y2) {};
\node[circ] at (x10y4) {};
\node[circ] at (x10y6) {};
\node[circ] at (x9y6) {};
\node[circ] at (x7y2) {};
\node[anchor=south west, inner sep=2pt, yshift=4pt, circuitnotemark, font=\LARGE\bfseries] (circuittitle) at (current bounding box.north west) {X};
\path (circuittitle.south west) rectangle ++(15.453,0.853);
\node[anchor=north east, inner sep=2pt, circuitnotemark, font=\large] (circuitstamp) at ([yshift=-0.593cm]current bounding box.south east) {X};
\path (circuitstamp.north east) rectangle ++(-5.08,-0.593);
\end{circuitikz}
\end{document}
